\begin{center}\begin{tikzpicture}[font=\small,>=stealth,line join=round]
\draw[gray!40,rounded corners=3pt] (0,0) rectangle (14.2,8.8);\node at (7.1,8.200000000000001) {补充证明：同一底数，正负指数的函数值互为倒数};
\node at (3.1500000000000004,7.3) {$y=x^2$};\begin{scope}[shift={(1.3533333333333335,3.3000000000000003)},x=1.6333333333333335cm,y=0.6cm]\draw[->,gray] (-0.4,0)--(2.6,0) node[right]{$x$};\draw[->,gray] (0,-0.5)--(0,5) node[above]{$y$};\draw[blue,thick,solid,domain=0.1:2.2,samples=180] plot(\x,{(\x)^2});\fill[orange!85!black] (0.5,0.25) circle(2pt);\fill[orange!85!black] (2,4) circle(2pt);\node[below=4pt] at (0.5,-0.3){$\frac12$};\node[below=4pt] at (2,-0.3){$2$};\end{scope}\node[align=center,text width=5.8cm] at (3.1500000000000004,1.7) {$x=\frac12$：$y=\frac14$\\$x=2$：$y=4$};\node at (9.95,7.3) {$y=x^{-2}$};\begin{scope}[shift={(8.153333333333334,3.3000000000000003)},x=1.6333333333333335cm,y=0.6cm]\draw[->,gray] (-0.4,0)--(2.6,0) node[right]{$x$};\draw[->,gray] (0,-0.5)--(0,5) node[above]{$y$};\draw[orange!85!black,thick,solid,domain=0.46:2.3,samples=180] plot(\x,{1/((\x)^2)});\fill[orange!85!black] (0.5,4) circle(2pt);\fill[orange!85!black] (2,0.25) circle(2pt);\node[below=4pt] at (0.5,-0.3){$\frac12$};\node[below=4pt] at (2,-0.3){$2$};\end{scope}\node[align=center,text width=5.8cm] at (9.95,1.7) {$x=\frac12$：$y=4$\\$x=2$：$y=\frac14$};
\node[text width=12.6cm,align=center] at (7.1,.65) {两图同尺度：同一横坐标处，两高度相乘恒为1。一般恒等式是$x^{-\alpha}=1/x^\alpha$。};
\end{tikzpicture}\end{center}
